\section{Reproducibility}

All results in this paper are based on explicit finite computations on
combinatorial data.

\subsection{Data}

The constructions use:
\begin{itemize}
\item a fixed labeling of the dodecahedron (Appendix~C),
\item an explicit edge list for $\Gfifteen$ (Appendix~A),
\item a signing $\sigfun$ defining the cocycle (Appendix~A).
\end{itemize}

\subsection{Deterministic procedures}

All computations are deterministic and involve:
\begin{itemize}
\item graph construction from flag transport,
\item quotient identification,
\item switching operations on signings,
\item exhaustive enumeration over $2^{14}$ switching configurations,
\item matrix computations on integer-valued data.
\end{itemize}

No randomness or probabilistic methods are used.

\subsection{Verification steps}

The main verifications can be reproduced by:
\begin{itemize}
\item constructing $\Gsixty$ from flag data and checking degree and connectivity,
\item computing the quotients $\Gsixty \to \Gthirty \to \Gfifteen$,
\item verifying $\Gfifteen \cong L(\mathrm{Petersen})$,
\item computing the signed lift from $\sigfun$,
\item enumerating switching configurations to confirm minimal support $6$,
\item computing the matrices $A$ and $M$ and checking
\[
MM^{\mathsf T} = A^3 + 2A^2 + 2I.
\]
\end{itemize}

\subsection{Implementation}

All computations are implemented in Python using explicit adjacency
lists and incidence structures. The scripts operate directly on the
data described in the appendices and require no external libraries
beyond standard numerical and combinatorial functionality.

\subsection{Scope}

The results are fully reproducible from the data and procedures given
in this paper. Any implementation following the same definitions will
produce identical outputs.
